% Appendix A -- the five-step decision logic of the capstone as a flowchart.
% Two dashed cuts mark the switches that turn the hybrid into the local-only and
% the replan-only strategy of the experiment matrix; the dotted rail is the path
% that the replan-only strategy takes once layer 3 is switched off.
\begin{tikzpicture}[x=1cm,y=1cm,font=\footnotesize,
  logicstep/.style={rounded corners=3pt,thick,align=center,inner sep=5pt,
               font={\scriptsize\sffamily\hyphenpenalty=10000\exhyphenpenalty=10000},
               text width=4.2cm},
  test/.style={diamond,aspect=3,draw=sbOrange,fill=sbLightOrange,thick,align=center,
               inner sep=1pt,
               font={\scriptsize\sffamily\hyphenpenalty=10000\exhyphenpenalty=10000},
               text width=2.5cm},
  trig/.style={font={\scriptsize\hyphenpenalty=10000\exhyphenpenalty=10000},
               text=sbGray,align=left,text width=3.1cm},
  lab/.style={font=\scriptsize,text=sbGray,inner sep=2pt,fill=white},
  cut/.style={draw=sbRed,very thick,dashed},
  cutlab/.style={font={\scriptsize\sffamily\hyphenpenalty=10000\exhyphenpenalty=10000},
               text=sbRed,align=center,inner sep=2pt}]

\node[logicstep,draw=sbBlue,fill=sbLightBlue]    (s1) at (0, 5.60) {\textbf{1} Follow the nominal \cbs/\ecbs path};
\node[test]                                      (d1) at (0, 3.93) {risk inside the safety horizon?};
\node[logicstep,draw=sbOrange,fill=sbLightOrange](s2) at (0, 1.76) {\textbf{2} \orca/DWA: local avoidance velocity};
\node[logicstep,draw=sbOrange,fill=sbLightOrange](s3) at (0, 0.35) {\textbf{3} Try to reconnect to the nominal path};
\node[test,draw=sbGreen,fill=sbLightGreen]       (d2) at (0,-1.32) {reconnection acceptable?};
\node[logicstep,draw=sbGreen,fill=sbLightGreen]  (s4) at (0,-3.49) {\textbf{4} \dstarlite/\astar replan from the current state};
\node[logicstep,draw=sbBlue,fill=sbLightBlue]    (s5) at (0,-4.90) {\textbf{5} Re-check inter-agent conflicts};

\draw[sbflow] (s1) -- (d1);
\draw[sbflow] (d1) -- node[lab,right,pos=0.30] {yes} (s2);
\draw[sbflow] (s2) -- (s3);
\draw[sbflow] (s3) -- (d2);
\draw[sbflow] (d2) -- node[lab,right,pos=0.30] {no} (s4);
\draw[sbflow] (s4) -- (s5);

% "no risk" and "reconnected" both return to step 1 along the right-hand rail
\draw[sbflow,rounded corners=4pt] (d1.east) -- node[lab,above,pos=0.25] {no} (4.6,3.93) -- (4.6,5.6) -- (s1.east);
\draw[sbflow,rounded corners=4pt] (d2.east) -- node[lab,above,pos=0.2] {yes} (5.6,-1.32) -- (5.6,6.1) -- (s1.north);
\draw[sbflow,rounded corners=4pt] (s5.east) -- (6.6,-4.90) -- (6.6,6.5) -- (0,6.5) -- (s1.north);

% the three triggers, on the left
\node[trig,anchor=east] at (-3.0,3.93)
  {\textbf{Trigger 1}\\ the predicted time to collision falls below the safety horizon $\ttc$};
\node[trig,anchor=east] at (-3.0,-1.32)
  {\textbf{Trigger 2}\\ reconnection cost above a bound\\[2pt]
   \textbf{Trigger 3}\\ formation or communication constraint violated};

% the two switches of the "avoidance strategy" axis
\draw[cut] (-1.05,2.82) -- (1.05,2.61);
\node[cutlab,anchor=west] at (1.2,2.72) {replan-only:\\ layer 3 off};
\draw[cut] (-1.05,-2.43) -- (1.05,-2.64);
\node[cutlab,anchor=west] at (1.2,-2.53) {local-only:\\ layer 4 off};

% with layer 3 off the yes branch runs straight to the replanner
\draw[sbflow,dotted,draw=sbRed,rounded corners=4pt]
  (0,3.06) -- (-2.7,3.06) -- (-2.7,-3.49) -- (s4.west);
\node[cutlab,rotate=90] at (-2.95,0.9) {replan-only};
\end{tikzpicture}
